%lezione 12 del 10 aprile 2026
%-----------------------------------------------------------------------------------------------------
% vd lez 11 per vero inizio di questa lezione, appena prima dela "diagonalizzazione" di X (line 133).
\subsection{Determinanti funzionali}
Come visto dalla relazione (\ref{eq:K_det}), un altro modo di stimare l'approssimazione semiclassica
di propagatori quantistici è trovare il determinante della forma quadratica $X$ associata a $S_2[z]$.

Si prende sempre come esempio l'oscillatore armonico, con
$$X = \dv[2]{}{t}+\omega^2$$
Per evitare difficoltà con la fase, si effettua una rotazione di Wick $t=-i\tau$, che manda
$$X\to X_e = - \dv[2]{}{\tau} + \omega^2$$
Dunque l'integrale della sezione precedente diventa
$$I = \int_{}^{}\mathcal{D}z \exp\left\{-\frac{m}{2\hbar}\int_{0}^{T_e}d\tau z
\left(-\dv[2]{}{\tau}+\omega^2\right)z\right\}$$
Come detto in precedenza, $\det(X_e) = \prod_{}^{}\lambda(X_e)$ e per $X_e$ vale che
$$X_e \sin\frac{\pi n\tau}{T_e}=\left(\frac{\pi^2n^2}{T_e^2} +\omega^2\right) \sin\frac{\pi n\tau}{T_e}$$
Il determinante è dunque
$$\det(X_e) = \prod_{n=1}^{\infty}\left(\frac{\pi^2n^2}{T_e^2} + \omega^2\right) =
\prod_{n=1}^{\infty}\left(\frac{T^2_e}{\pi^2n^2}\right)\prod_{n=1}^{\infty}
\left(1+\frac{\omega^2T_e^2}{\pi^2n^2}\right)$$
\textit{Oss.} Con questa visione è chiaro che $\det(-\dv[2]{}{\tau}) = \prod_{n=1}^{\infty}\left(
\frac{T^2_e}{\pi^2n^2}\right)$: questo è il determinante funzionale della forma quadratica associata
alla particella libera.

Tornando dalla rotazione, $T_e = iT$ e il calcolo delle produttorie è il medesimo della sezione
precedente.
\subsubsection{Regolarizzazione con funzione zeta}
In questo e molti altri problemi di fisica teorica, ci si imbatte in oggetti matematici non convergenti
e, per trovare un significato al loro risultato, è necessario regolarizzarli.\\
Si prenda come esempio l'appena citato determinante funzionale per la particella libera
$$\det(X) = \prod_{n=1}^{\infty}\left(\frac{T^2}{\pi^2n^2}\right)$$
questa produttoria è divergente, ma, sfruttando la continuazione analitica della $\zeta$ di Riemann,
è possibile associare a essa un risultato finito.\\
Detti $\lambda_n$ gli autovalori di $X$ si consideri la funzione
$$\zeta_X(s):=\sum_{n=1}^{\infty}\frac{1}{\lambda^s_{n}}$$
La sua derivata è
$$\zeta_X'(s) = \dv[]{}{s}\left(\sum_{n=1}^{\infty}e^{-s\ln(\lambda_n)}\right)=
\sum_{n=1}^{\infty}-\frac{1}{\lambda_n^2}\ln(\lambda_n) \rightarrow \zeta_X'(0) =
-\sum_{n=1}^{\infty}\ln(\lambda_n)$$
e dunque $\det(X) = e^{-\zeta_X'(0)}$

Sosituendo gli autovalori $\lambda_n=\frac{T^2}{\pi^2n^2}$ si trova
$$\zeta_X(s) = \left(\frac{T}{\pi}\right)^{2s}\sum_{n=1}^{\infty}\frac{1}{n^{2s}} =
\left(\frac{T}{\pi}\right)^{2s}\zeta(2s)$$
La sua derivata in zero vale
$$\zeta_X'(0) =\eval{ \left(\frac{T}{\pi}\right)^{2s}\left[2\ln\left(\frac{T}{\pi}\right)\zeta(2s)+2\zeta'(2s)
\right]}_{s=0} = 2\ln\left(\frac{T}{\pi}\right)\zeta(0) + 2\zeta'(0)$$
ricordando che, per continuazione analitica, $\zeta(0) = -\frac{1}{2}$ e $\zeta'(0) = \ln(2\pi)\zeta(0)$,
si trova infine
$$\zeta_X'(0) = -\ln(2T)\rightarrow\det(X) = 2T$$
che è il risultato atteso a meno di costanti di normalizzazione.
\section{Rotazione di Wick e meccanica statistica}

